\section{Use Cases}

Apache Kafka has become the de facto standard event backbone across a wide range of industries.
Table~\ref{tab:kafka-usecases} summarises the canonical use cases; the subsections below explain
each pattern in more detail.

\begin{table}[H]
  \centering
  \caption{Common Apache Kafka use cases.}
  \label{tab:kafka-usecases}
  \begin{tabularx}{\textwidth}{lX}
    \toprule
    \textbf{Use Case}              & \textbf{Description} \\
    \midrule
    Real-time event streaming      & Collect click streams, IoT telemetry, or application logs and process them with Flink or Spark. \\
    Microservice messaging         & Decouple services with an async pub/sub channel instead of synchronous REST or gRPC calls. \\
    Change Data Capture (CDC)      & Stream database change events (via Debezium) to keep downstream systems in sync. \\
    Log aggregation                & Centralise logs from hundreds of servers into a single, queryable stream. \\
    Metrics \& monitoring          & Route time-series metrics to alerting pipelines and observability stores. \\
    Event sourcing                 & Use a compacted topic as the system of record; reconstruct application state by replaying events. \\
    Stream ETL                     & Read raw events, transform and enrich them in flight (Flink/Spark), and land clean data in a warehouse. \\
    \bottomrule
  \end{tabularx}
\end{table}

\subsection{Real-time Event Streaming}

High-velocity event streams — user click sequences on an e-commerce site, GPS pings from a
fleet of vehicles, sensor readings from industrial IoT devices — are a natural fit for Kafka.
Events are published by producers (edge devices, application servers, SDKs) and consumed by one
or more independent processing engines (Flink, Spark, or custom consumers) with millisecond ingestion latency.
A key advantage over direct HTTP calls is \textbf{durability}: if the downstream processor
restarts or falls behind, it replays events from its last committed offset rather than losing them.
This pattern is the one used in Streamify, where Eventsim publishes music events to three Kafka
topics and PyFlink consumes them for real-time processing.

\subsection{Microservice Messaging}

In a synchronous microservice architecture, service A calls service B's REST API directly.
When service B is slow or unavailable, service A blocks or drops the request.
Kafka replaces this point-to-point call with a shared topic: service A publishes an event and
returns immediately; service B (and any other subscriber) processes it asynchronously at its own pace.
This \textbf{temporal decoupling} eliminates cascading failures and allows services to be deployed,
scaled, or replaced independently.
Unlike traditional message queues (RabbitMQ, ActiveMQ), Kafka retains events after delivery,
enabling late consumers to join the system and catch up on historical events.

\subsection{Change Data Capture (CDC)}

CDC is the practice of streaming every INSERT, UPDATE, and DELETE from a relational database
into downstream systems in near real time.
\textbf{Debezium}, the leading open-source CDC tool, connects to the database's replication log
(PostgreSQL WAL, MySQL binlog, Oracle LogMiner) and publishes change events to Kafka topics.
Consumers can then build materialised views, search indexes, or analytical aggregates that
stay current within seconds of the source database change — without polling the database.
Kafka's log compaction feature is particularly valuable here: a compacted topic retains the
latest record per key, effectively functioning as a full replica of the current database state.

\subsection{Log Aggregation}

Distributed systems running hundreds of servers produce logs in disparate formats across many
machines.
Shipping these logs directly to a central store (Elasticsearch, Splunk, S3) creates tight
coupling between log producers and the storage backend.
Kafka acts as a \textbf{log buffer}: lightweight agents (Filebeat, Fluentd) tail log files and
publish entries to a Kafka topic; consumers (Logstash, Flink) parse, filter, and route them to
storage.
This decoupling means the log storage backend can be changed or scaled without touching the
hundreds of log-producing services.

\subsection{Metrics and Monitoring}

Time-series metrics (CPU utilisation, request latency percentiles, error rates) are naturally
event-shaped data.
Publishing metrics to Kafka allows them to fan out to multiple consumers simultaneously:
a Prometheus-compatible consumer for real-time alerting, a Flink job for anomaly detection,
and an S3 consumer for long-term retention — all from a single Kafka topic, with no additional
load on the metrics producer.
Kafka's configurable retention window also enables retrospective analysis: if an alert fires at
3\,am, the raw metric stream can be replayed from before the incident began.

\subsection{Event Sourcing}

In an event-sourced system, the authoritative state of an entity is not stored as a mutable
database row but as an ordered sequence of immutable events in Kafka (e.g.\ \texttt{OrderPlaced},
\texttt{PaymentReceived}, \texttt{OrderShipped}).
Application state is reconstructed on demand by replaying the event log.
Kafka's \textbf{log compaction} mode retains the latest record for each key, providing an
efficient snapshot of the current state without replaying the full history.
This pattern provides a natural audit log, simplifies debugging (replay past events to reproduce
bugs), and enables event-driven projections to be rebuilt at any time.

\subsection{Stream ETL}

Traditional ETL (Extract → Transform → Load) runs nightly batch jobs that move data from
operational systems into data warehouses.
Stream ETL (more accurately described as ELT) inverts this: events land in Kafka continuously,
a stream processor (Flink or Spark) transforms and enriches them in real time, and clean records
are written to the warehouse within seconds of the original event.
This is the pattern used in Streamify: Eventsim → Kafka → PyFlink (enrich, deduplicate) →
Snowflake (raw staging) → dbt (star schema) → Power BI.
The result is a BI dashboard whose data is fresher than any nightly batch pipeline could achieve.
